\documentclass[11pt,oneside]{book}
% Shared preamble for the Carina theory manual.
%
% Every chapter writes tensors and matrices through the macros defined here,
% so the notation is set in one place and cannot drift between chapters.

\usepackage[T1]{fontenc}
\usepackage{lmodern}
\usepackage{amsmath, amssymb, amsthm}
\usepackage{boldtensors}
\usepackage{booktabs}
\usepackage{microtype}
% `algorithmic` rather than `algpseudocode`: the latter lives in algorithmicx,
% which is not in a base TeX Live installation.  This manual builds against
% texlive-latex-recommended with no extra package installs.
\usepackage{algorithm}
\usepackage{algorithmic}
\usepackage[margin=1.15in]{geometry}
\usepackage[hidelinks]{hyperref}

% File paths set in typewriter cannot hyphenate, so a paragraph containing one
% can overflow the measure.  Let TeX stretch the interword space as a last
% resort rather than run into the margin.
\setlength{\emergencystretch}{2.5em}

% --------------------------------------------------------------------------
% Notation
% --------------------------------------------------------------------------
% boldtensors makes \bf-prefixed Greek and upright symbols bold, which is what
% carries the tensor/matrix distinction throughout.  Second-order tensors and
% matrices are bold upright capitals, vectors bold lowercase, fourth-order
% tensors blackboard bold.

% boldtensors makes `~` and `"` active in math mode: `~X` sets X in the bold
% tensor font, `"X` in blackboard bold.  That is the package's own interface;
% wrapping it in named macros keeps the active characters out of the chapter
% sources, where they would be easy to mistype and hard to grep for.

\newcommand{\tensor}[1]{~#1}                  % 2nd-order tensor / matrix
\newcommand{\vect}[1]{~#1}                    % vector
\newcommand{\fourth}[1]{"#1}                  % 4th-order tensor

% Continuum quantities
\newcommand{\Fdef}{\tensor{F}}                % deformation gradient
\newcommand{\Piola}{\tensor{P}}               % first Piola--Kirchhoff stress
\newcommand{\Cauchy}{\tensor{\sigma}}         % Cauchy stress
\newcommand{\Kirch}{\tensor{\tau}}            % Kirchhoff stress
\newcommand{\strain}{\tensor{\varepsilon}}    % small strain
\newcommand{\Cright}{\tensor{C}}              % right Cauchy--Green
\newcommand{\Ident}{\tensor{I}}               % identity tensor
\newcommand{\disp}{\vect{u}}                  % displacement field
\newcommand{\Grad}{\nabla}

% Discrete quantities
\newcommand{\Kmat}{\tensor{K}}                % tangent stiffness
\newcommand{\Mmat}{\tensor{M}}                % mass matrix
\newcommand{\Keff}{\tensor{K}_{\mathrm{eff}}} % Newmark effective stiffness
\newcommand{\Rvec}{\vect{R}}                  % residual
\newcommand{\Uvec}{\vect{u}}                  % discrete displacement
\newcommand{\Vvec}{\vect{v}}                  % discrete velocity
\newcommand{\Avec}{\vect{a}}                  % discrete acceleration
\newcommand{\Fint}{\vect{f}^{\mathrm{int}}}
\newcommand{\Fext}{\vect{f}^{\mathrm{ext}}}
\newcommand{\Pmat}{\tensor{P}}                % prolongation
\newcommand{\Rrest}{\tensor{R}}               % restriction
\newcommand{\Dmat}{\tensor{D}}                % diagonal of an operator
\newcommand{\Bnull}{\tensor{B}}               % near-nullspace
\newcommand{\Prec}{\tensor{M}}                % preconditioner

% Operators
\DeclareMathOperator{\dev}{dev}
\DeclareMathOperator{\tr}{tr}
\DeclareMathOperator{\diag}{diag}
\DeclareMathOperator*{\argmin}{arg\,min}
\newcommand{\norm}[1]{\left\lVert #1 \right\rVert}
\newcommand{\abs}[1]{\left\lvert #1 \right\rvert}
\newcommand{\inner}[2]{\left\langle #1, #2 \right\rangle}
\newcommand{\dbldot}{\mathbin{:}}

% Cross-references into the implementation.  Every chapter names the files that
% implement it, so the two can be checked against each other.
%
% Set in the text flow rather than as a \marginpar: the margins here are too
% narrow for a path like `src/nonlinear_solvers.jl`, and a marginpar that does
% not fit is silently clipped at the page edge rather than reported.
\newcommand{\impl}[1]{%
  \vspace{-0.6em}%
  {\raggedleft\small\textit{Implementation:}~\texttt{#1}\par}%
  \vspace{0.4em}%
}
\newcommand{\code}[1]{\texttt{#1}}

\theoremstyle{definition}
\newtheorem{remark}{Remark}[chapter]
\theoremstyle{plain}
\newtheorem{proposition}{Proposition}[chapter]


\title{\textbf{Carina}\\[0.4em]
       \large Theory Manual\\[0.3em]
       \normalsize Formulation, constitutive models, and solution algorithms}
\author{Alejandro Mota}
\date{\today}

\begin{document}
\maketitle

\frontmatter
\chapter*{About this manual}
\addcontentsline{toc}{chapter}{About this manual}

Carina's documentation is split by the question it answers.  The user manual
(\code{docs/}, built with Documenter) answers \emph{what do I write in the input
deck}.  The benchmark report (\code{benchmark\_report.md}) answers \emph{how
fast is it, and on what hardware}.  This manual answers \emph{why the
algorithms work}.

It is written to be read alongside the source.  Each algorithm is stated in the
form the code actually implements --- not the form most convenient to typeset
--- and the margin of each section names the file implementing it.  Where the
implementation departs from a textbook statement, the departure is the point of
the section rather than a footnote to it: those departures are usually the
result of a measurement, and the measurement is cited.

Nothing here is a substitute for reading the code, and the code is the
authority wherever the two disagree.  If you find such a disagreement, it is a
defect in this manual.

\tableofcontents

\mainmatter

\chapter{Notation}
\label{ch:notation}

Bold upright capitals denote second-order tensors and matrices, bold lowercase
denotes vectors, and blackboard bold denotes fourth-order tensors.  Bold Greek
follows the same convention through the \code{boldtensors} package: $\Cauchy$
is the Cauchy stress tensor, $\sigma$ a scalar.

\section*{Continuum quantities}

\begin{tabular}{@{}ll@{}}
\toprule
Symbol & Meaning \\
\midrule
$\vect{X}$, $\vect{x}$        & material and spatial position \\
$\disp$                       & displacement field, $\vect{x} = \vect{X} + \disp$ \\
$\Fdef = \Ident + \Grad\disp$ & deformation gradient \\
$J = \det\Fdef$               & Jacobian of the deformation \\
$\Cright = \Fdef^{\mathsf{T}}\Fdef$ & right Cauchy--Green tensor \\
$\Piola$                      & first Piola--Kirchhoff stress \\
$\Cauchy$, $\Kirch = J\Cauchy$ & Cauchy and Kirchhoff stress \\
$\Psi$                        & strain-energy density per unit reference volume \\
$\fourth{A} = \partial\Piola/\partial\Grad\disp$ & material tangent \\
$\rho_0$                      & reference mass density \\
\bottomrule
\end{tabular}

\section*{Discrete quantities}

\begin{tabular}{@{}ll@{}}
\toprule
Symbol & Meaning \\
\midrule
$n$                     & number of unknown (free) degrees of freedom \\
$\Uvec, \Vvec, \Avec$   & nodal displacement, velocity, acceleration \\
$\Rvec$                 & residual (out-of-balance force) \\
$\Kmat = \partial\Rvec/\partial\Uvec$ & tangent stiffness \\
$\Mmat$                 & mass matrix \\
$\Keff$                 & effective operator of an implicit step \\
$c_M$                   & mass coefficient of $\Keff$, Section~\ref{sec:newmark} \\
$\Dmat = \diag(\Kmat)$  & diagonal of an operator \\
$\Prec$                 & preconditioner, $\Prec \approx \Kmat^{-1}$ \\
$\Pmat$, $\Rrest$       & prolongation and restriction \\
$\Bnull$                & near-nullspace, Section~\ref{sec:nullspace} \\
$\eta_k$                & forcing term at Newton iteration $k$, Chapter~\ref{ch:inexact} \\
$\lambda_{\max}$        & largest eigenvalue of $\Dmat^{-1}\Kmat$ \\
\bottomrule
\end{tabular}

\section*{Index conventions}

Subscript $k$ indexes nonlinear (Newton) iterations, subscript $i$ or $j$
indexes linear (Krylov) iterations, subscript $n$ indexes time steps, and
superscript $\ell$ indexes multigrid levels.  A quantity written without a time
subscript is understood at $t_{n+1}$, the step being solved.

Norms are Euclidean unless a subscript says otherwise; $\norm{\vect{r}}_{\Prec}
= \sqrt{\vect{r}^{\mathsf{T}}\Prec\vect{r}}$ denotes the norm induced by a
preconditioner, which is the quantity preconditioned conjugate gradients
actually monitors (Section~\ref{sec:pcg}).

\section*{Free and constrained degrees of freedom}

Carina imposes Dirichlet conditions by \emph{elimination}, never by penalty.
The global degree-of-freedom set partitions into unknowns and prescribed
values, and every operator named above is understood as the unknown--unknown
block unless stated otherwise.  Chapter~\ref{ch:formulation} makes this precise
and Section~\ref{sec:elimination} explains why the choice matters for every
Krylov method in this manual.

\chapter{Formulation}
\label{ch:formulation}
\impl{src/physics.jl}

\section{Strong and weak form}

Carina solves finite-deformation solid mechanics in the reference
configuration.  Linear momentum balance reads
\begin{equation}
  \Grad \cdot \Piola + \rho_0 \vect{b} = \rho_0 \ddot{\disp},
  \label{eq:momentum}
\end{equation}
with $\Piola$ the first Piola--Kirchhoff stress, $\vect{b}$ a body force per
unit mass, and $\rho_0$ the reference density.  Multiplying by an admissible
variation $\delta\disp$ vanishing on the Dirichlet boundary $\Gamma_u$ and
integrating by parts gives the weak form: find $\disp$ with $\disp =
\bar{\disp}$ on $\Gamma_u$ such that
\begin{equation}
  \int_{\Omega_0} \Piola \dbldot \Grad \delta\disp \, \mathrm{d}V
  + \int_{\Omega_0} \rho_0\, \ddot{\disp} \cdot \delta\disp \, \mathrm{d}V
  = \int_{\Omega_0} \rho_0 \vect{b}\cdot\delta\disp \, \mathrm{d}V
  + \int_{\Gamma_t} \bar{\vect{t}} \cdot \delta\disp \, \mathrm{d}A
  \label{eq:weak}
\end{equation}
for all admissible $\delta\disp$, where $\bar{\vect{t}}$ is the prescribed
traction on $\Gamma_t$.

The kinematics enter only through the deformation gradient
\begin{equation}
  \Fdef = \Ident + \Grad\disp ,
\end{equation}
and the constitutive response enters only through $\Piola(\Fdef)$.  This is the
entire coupling surface between Chapter~\ref{ch:constitutive} and everything
downstream: a constitutive model supplies $\Piola$ given $\Grad\disp$, and
optionally its derivative.

\section{Discretization}

Introducing isoparametric shape functions $N_a$ over an element mesh and
writing $\disp \approx \sum_a N_a \Uvec_a$, \eqref{eq:weak} becomes the
semi-discrete system
\begin{equation}
  \Mmat \Avec + \Fint(\Uvec) = \Fext(t),
  \label{eq:semidiscrete}
\end{equation}
with
\begin{equation}
  \Mmat_{ab} = \int_{\Omega_0} \rho_0 N_a N_b \, \mathrm{d}V,
  \qquad
  \Fint_a(\Uvec) = \int_{\Omega_0} \Piola(\Fdef) \cdot \Grad N_a \, \mathrm{d}V .
\end{equation}
Both integrals are evaluated by Gauss quadrature element by element.  The
residual that every solver in this manual drives to zero is
\begin{equation}
  \Rvec(\Uvec) = \Fext - \Fint(\Uvec) - \Mmat\Avec ,
  \label{eq:residual}
\end{equation}
with the inertial term absent in the quasi-static case
(Section~\ref{sec:quasistatic}).

\begin{remark}
Carina stores $\Rvec$ already negated relative to the force residual, so the
Newton system is $\Kmat\,\Delta\Uvec = \Rvec$ with a plus sign throughout the
implementation.  The sign convention is stated in \code{linear\_solvers.jl} at
the dispatch point and is followed uniformly here.
\end{remark}

\section{Consistent tangent}

Differentiating \eqref{eq:residual} with respect to the nodal displacements
gives the tangent stiffness
\begin{equation}
  \Kmat_{ab} = \frac{\partial \Fint_a}{\partial \Uvec_b}
             = \int_{\Omega_0} \Grad N_a \cdot \fourth{A} \cdot \Grad N_b \,
               \mathrm{d}V,
  \qquad
  \fourth{A} = \frac{\partial \Piola}{\partial \Grad\disp} .
  \label{eq:tangent}
\end{equation}
The fourth-order material tangent $\fourth{A}$ is the only object a
constitutive model must supply beyond the stress itself; how Carina obtains it
--- in closed form, by automatic differentiation, or by finite differences ---
is the subject of Section~\ref{sec:tangent-strategies}.

\section{Dirichlet conditions by elimination}
\label{sec:elimination}

Prescribed displacements are imposed by removing the corresponding equations
and unknowns, not by adding a penalty term or a Lagrange multiplier.
Partitioning the degrees of freedom into unknowns $\mathrm{f}$ and prescribed
values $\mathrm{p}$,
\begin{equation}
  \begin{bmatrix}
    \Kmat_{\mathrm{ff}} & \Kmat_{\mathrm{fp}} \\
    \Kmat_{\mathrm{pf}} & \Kmat_{\mathrm{pp}}
  \end{bmatrix}
  \begin{bmatrix} \Delta\Uvec_{\mathrm{f}} \\ \Delta\Uvec_{\mathrm{p}} \end{bmatrix}
  =
  \begin{bmatrix} \Rvec_{\mathrm{f}} \\ \Rvec_{\mathrm{p}} \end{bmatrix},
\end{equation}
only the first block row is solved, with $\Delta\Uvec_{\mathrm{p}}$ known.  The
cross term $\Kmat_{\mathrm{fp}}\Delta\Uvec_{\mathrm{p}}$ is carried
automatically because the element kernels read displacement fields that already
hold the prescribed values at constrained nodes; nothing assembles
$\Kmat_{\mathrm{fp}}$ explicitly.

This choice is not cosmetic.  A penalty imposition multiplies selected
diagonal entries by a large factor --- $10^{15}$ is typical --- which raises
the condition number of the operator by the same factor and destroys the
convergence rate of every Krylov method in Chapter~\ref{ch:krylov}.  The
eliminated system carries no artificial conditioning at all, so the iteration
counts reported throughout this manual reflect the physics of the problem
rather than the bookkeeping of the boundary conditions.

\begin{remark}
One place in Carina does see penalty-scale entries: the assembled Chebyshev
path, which forms $\Kmat$ including constrained rows before scaling.  That is
why the polynomial there is built on the symmetrically scaled operator
$\Dmat^{-1/2}\Kmat\Dmat^{-1/2}$, whose eigenvalues are $O(1)$ even when
$\Kmat$'s are not.  See Section~\ref{sec:chebyshev}.
\end{remark}

\section{Mass matrix and lumping}

The consistent mass of \eqref{eq:semidiscrete} is used by the implicit dynamic
integrator.  The explicit integrator instead uses the row-sum lumped mass
\begin{equation}
  \Mmat^{\mathrm{L}}_{aa} = \sum_b \Mmat_{ab},
\end{equation}
which is diagonal and therefore invertible in $O(n)$ with no linear solve at
all --- the property that makes explicit integration attractive
(Section~\ref{sec:central-difference}).

\chapter{Constitutive models}
\label{ch:constitutive}
\impl{src/physics.jl, src/materials.jl}

Carina delegates constitutive evaluation to \code{ConstitutiveModels.jl}.  The
interface is narrow by design: a model returns $\Piola$ given $\Grad\disp$ and,
for models carrying internal variables, an updated state.  Four models are
wired into the input parser.

\section{Linear elastic}

Small-strain isotropic elasticity, $\Cauchy = \lambda\tr(\strain)\Ident +
2\mu\strain$ with $\strain = \tfrac12(\Grad\disp +
\Grad\disp^{\mathsf{T}})$.  Its tangent is constant, which Carina exploits: for
a linear-elastic model the stiffness is assembled once and the factorization
--- or, on GPU, the Cholesky factor uploaded for triangular solves
(Section~\ref{sec:gpu-cholesky}) --- is cached for the entire run.

\section{Neo-Hookean}

The compressible neo-Hookean strain energy
\begin{equation}
  \Psi(\Fdef) = \frac{\mu}{2}\left(\tr\Cright - 3\right) - \mu \ln J
              + \frac{\lambda}{2}\left(\ln J\right)^2
\end{equation}
gives
\begin{equation}
  \Piola = \mu\left(\Fdef - \Fdef^{-\mathsf{T}}\right)
         + \lambda \ln J \, \Fdef^{-\mathsf{T}} .
\end{equation}
This is the workhorse for the benchmark problems throughout the performance
campaigns, and the one model whose tangent \emph{action} Carina writes in
closed form (Section~\ref{sec:tangent-strategies}).

\section{Hencky}

Logarithmic (true) strain elasticity, $\Kirch = \lambda\tr(\vect{h})\Ident +
2\mu\vect{h}$ with $\vect{h} = \tfrac12\ln\Cright$.  The matrix logarithm is
evaluated by inverse scaling-and-squaring with Pad\'e approximants rather than
by spectral decomposition: the spectral route loses accuracy and
differentiability when eigenvalues of $\Cright$ coalesce, which is exactly the
neighbourhood of the undeformed state.

\section{Finite-deformation \texorpdfstring{$J_2$}{J2} plasticity}

Rate-independent $J_2$ flow with linear isotropic hardening, integrated by a
radial-return map at each quadrature point.  The trial elastic state is
computed, and if the yield function
\begin{equation}
  f = \norm{\dev \Kirch^{\mathrm{tr}}} - \sqrt{\tfrac23}\left(\sigma_y + H\,\bar{\varepsilon}^p\right)
\end{equation}
is positive the plastic multiplier is solved for and the stress projected back
onto the yield surface.  The model carries internal state --- the plastic
strain tensor and the equivalent plastic strain --- which makes it
\emph{path dependent}: the same displacement reached by different increments
gives different stresses.  This has two consequences that surface elsewhere in
this manual.  Large load steps miss the yield surface and must be prevented by
adaptive stepping (Section~\ref{sec:adaptive}); and the return map can raise
domain errors on non-physical trial states, which the solver treats as a step
failure rather than a crash.

\section{Tangent strategies}
\label{sec:tangent-strategies}

Equation~\eqref{eq:tangent} needs $\fourth{A} = \partial\Piola/\partial
\Grad\disp$ at every quadrature point.  Carina obtains it three different ways,
chosen by what the model can support.

\paragraph{Closed form.}  For neo-Hookean the directional derivative is
written analytically.  With $\delta\Fdef = \Grad\vect{v}$, and abbreviating
$a = \tr(\Fdef^{-1}\delta\Fdef)$ and $\tensor{W} =
\Fdef^{-\mathsf{T}}\delta\Fdef^{\mathsf{T}}\Fdef^{-\mathsf{T}}$,
\begin{equation}
  \delta\Piola = \mu\left(\delta\Fdef + \tensor{W}\right)
               + \lambda\, a\, \Fdef^{-\mathsf{T}}
               - \lambda \ln J\, \tensor{W} .
  \label{eq:nh-action}
\end{equation}
This form is what the matrix-free operator of Chapter~\ref{ch:matrixfree}
evaluates: it needs the tangent's action on \emph{one} direction, never the
full $81$-component $\fourth{A}$, so the closed form costs a handful of
$3\times3$ products instead of a rank-four object.

\paragraph{Automatic differentiation.}  For stateless hyperelastic models
generally, $\fourth{A}$ is obtained by differentiating \code{pk1\_stress}
through dual numbers.  This is exact to machine precision and requires nothing
from the model author.

\paragraph{Finite differences.}  Models carrying state --- $J_2$ plasticity ---
run a return map inside \code{pk1\_stress}, whose branches are not
differentiable in the ordinary sense.  Carina falls back to a forward finite
difference over the nine components of $\Grad\disp$,
\begin{equation}
  \fourth{A}_{ijkl} \approx
  \frac{\Piola_{ij}(\Grad\disp + h_{kl}\,\vect{e}_k\otimes\vect{e}_l) - \Piola_{ij}(\Grad\disp)}{h_{kl}},
\end{equation}
at a cost of nine stress evaluations per quadrature point.  The perturbation is
adaptive,
\begin{equation}
  h_{kl} = \max\left(\sqrt{\varepsilon_{\mathrm{mach}}}\,\abs{(\Grad\disp)_{kl}},\ h_{\min}\right),
  \qquad h_{\min} = 10^{-10},
\end{equation}
following the Sierra/SM convention.  The relative term dominates wherever the
strain is resolved; the floor prevents a vanishing perturbation --- and hence
catastrophic cancellation --- at components that happen to be near zero, which
in an undeformed region is most of them.

\begin{remark}
The finite-difference tangent is a \emph{tangent}, not the residual: it enters
only the Newton direction and the preconditioner, never the convergence test.
An approximate tangent slows Newton's asymptotic rate but cannot change the
solution it converges to, because that is fixed by $\Rvec = \vect{0}$.  The
same argument licenses the reduced-precision operator of
Section~\ref{sec:mixed-precision}.
\end{remark}

\chapter{Time integration}
\label{ch:integration}
\impl{src/integrators.jl}

Three integrators share one \code{predict!}\,/\,\code{evaluate!}\,/\,%
\code{correct!} protocol, so the nonlinear and linear machinery downstream is
written once.

\section{Quasi-static}
\label{sec:quasistatic}

Inertia is dropped from \eqref{eq:semidiscrete} and the load is applied in
pseudo-time increments:
\begin{equation}
  \Rvec(\Uvec) = \Fext(t_{n+1}) - \Fint(\Uvec) = \vect{0},
  \qquad \Kmat = \frac{\partial\Fint}{\partial\Uvec} .
\end{equation}
The operator is the bare tangent stiffness.  For a well-shaped mesh its
condition number grows like $(L/h)^2$ with the mesh refinement ratio, which is
what makes quasi-static problems the natural home for multigrid
(Chapter~\ref{ch:amg}) and the reason a diagonal preconditioner degrades so
visibly with problem size.

\section{Newmark-\texorpdfstring{$\beta$}{beta}}
\label{sec:newmark}

The Newmark family parameterizes the step by $\beta$ and $\gamma$:
\begin{align}
  \Uvec_{n+1} &= \Uvec_n + \Delta t\, \Vvec_n
                + \Delta t^2\left[\left(\tfrac12 - \beta\right)\Avec_n + \beta \Avec_{n+1}\right], \\
  \Vvec_{n+1} &= \Vvec_n + \Delta t\left[(1-\gamma)\Avec_n + \gamma \Avec_{n+1}\right].
\end{align}
Carina implements this in predictor--corrector form.  The predictor collects
everything known at $t_n$,
\begin{equation}
  \Uvec^{\mathrm{pred}} = \Uvec_n + \Delta t\, \Vvec_n
                          + \Delta t^2\left(\tfrac12 - \beta\right)\Avec_n,
  \qquad
  \Vvec^{\mathrm{pred}} = \Vvec_n + \Delta t (1-\gamma)\Avec_n,
\end{equation}
after which the acceleration follows from the displacement alone:
\begin{equation}
  \Avec_{n+1} = c_M\left(\Uvec_{n+1} - \Uvec^{\mathrm{pred}}\right),
  \qquad
  \boxed{\ c_M = \frac{1}{\beta \Delta t^2}\ }
  \label{eq:cM}
\end{equation}
and the corrector restores the velocity, $\Vvec_{n+1} = \Vvec^{\mathrm{pred}} +
\gamma\Delta t\,\Avec_{n+1}$.  Displacement is thus the only unknown, and the
Newton system for it has the \emph{effective operator}
\begin{equation}
  \Keff = \Kmat + c_M \Mmat .
  \label{eq:keff}
\end{equation}

Equation~\eqref{eq:cM} governs almost everything about how implicit dynamics
behaves numerically, and it is worth being explicit about why.  As $\Delta t$
shrinks, $c_M \sim \Delta t^{-2}$ grows without bound and $\Keff$ becomes
mass-dominated.  The mass matrix is well conditioned --- it is a Gram matrix of
shape functions, independent of the deformation --- so $\Keff$ inherits that
conditioning.  Three consequences recur throughout this manual:

\begin{itemize}
  \item A diagonal (Jacobi) preconditioner is unusually effective at small
        $\Delta t$, because $\diag(c_M\Mmat)$ captures the dominant term.
  \item Multigrid, whose value is mesh-independent iteration counts on a
        stiffness-dominated operator, has little left to remove; measurements
        confirm it does not repay its per-application cost in this regime
        (Section~\ref{sec:amg-when}).
  \item A low-rank quasi-Newton model of $\Keff^{-1}$ is viable at small
        $\Delta t$ and not at all in the quasi-static limit
        (Section~\ref{sec:lbfgs}).
\end{itemize}

\paragraph{Constrained degrees of freedom.}  The predictor is engineered so
that \eqref{eq:cM} reproduces the prescribed acceleration exactly at
constrained nodes: $\Uvec^{\mathrm{pred}}_{\mathrm{p}} = \Uvec_{\mathrm{p}} -
\ddot{g}(t_{n+1})/c_M$, whence $\Avec_{\mathrm{p}} = \ddot{g}(t_{n+1})$.  This
keeps the inertial cross-term $\Mmat_{\mathrm{fp}}\ddot{g}$ in the residual
without special-casing it, and it is why the integrator carries full-length
$\Uvec$, $\Vvec$, $\Avec$ arrays rather than free-DOF slices.

\paragraph{Stability.}  $\beta = \tfrac14$, $\gamma = \tfrac12$ (average
acceleration) is unconditionally stable and second-order accurate, and is the
default.  Numerical dissipation, when wanted, is available through the
HHT-$\alpha$ generalization the integrator carries.

\section{Central difference}
\label{sec:central-difference}

The explicit integrator advances
\begin{equation}
  \Uvec_{n+1} = \Uvec_n + \Delta t \Vvec_n + \tfrac12 \Delta t^2 \Avec_n,
  \qquad
  \Avec_{n+1} = \left(\Mmat^{\mathrm{L}}\right)^{-1}\Rvec(\Uvec_{n+1}),
\end{equation}
followed by $\Vvec_{n+1} = \Vvec_n + \Delta t\left[(1-\gamma)\Avec_n + \gamma
\Avec_{n+1}\right]$.  With a lumped mass the acceleration update is a
component-wise division: no linear solve, no nonlinear solve, no tangent.
Nothing in Chapters~\ref{ch:newton}--\ref{ch:amg} applies, which is why the
input parser ignores the \code{solver} block entirely for this integrator.

\paragraph{Critical time step.}  Stability is conditional on the CFL limit
\begin{equation}
  \Delta t \le \mathrm{CFL} \cdot \min_e \frac{\ell_e}{c_e},
  \qquad c_e = \sqrt{\frac{\lambda + 2\mu}{\rho_0}},
\end{equation}
with $\ell_e$ a characteristic element length and $c_e$ the dilatational wave
speed.  Carina recomputes this bound periodically during the run rather than
once at setup, because $\ell_e$ shrinks as elements compress.

\section{Adaptive time stepping}
\label{sec:adaptive}

Both implicit integrators shrink $\Delta t$ by a decrease factor when the
nonlinear solve fails and grow it by an increase factor, up to a ceiling, when
it succeeds.  A failure is any of: the status tree reporting
\textsc{Failed}~(Section~\ref{sec:status}), a non-finite residual, or a domain
error raised inside a constitutive model.

The mechanism matters most for $J_2$ plasticity, where it is not merely an
efficiency device but a correctness one.  Because the return map is path
dependent, a step large enough to traverse the yield surface in one increment
integrates the wrong path --- and converges perfectly well to the wrong answer.
Step-size control plus the line search of Section~\ref{sec:linesearch} is what
keeps the increments small enough for the return map to track the true path.

\chapter{Newton's method}
\label{ch:newton}
\impl{src/nonlinear\_solvers.jl}

\section{The iteration}

Newton--Raphson linearizes $\Rvec(\Uvec) = \vect{0}$ about the current iterate:
\begin{equation}
  \Kmat(\Uvec_k)\, \Delta\Uvec_k = \Rvec(\Uvec_k),
  \qquad
  \Uvec_{k+1} = \Uvec_k + \alpha_k \Delta\Uvec_k ,
  \label{eq:newton}
\end{equation}
with $\Kmat$ the tangent \eqref{eq:tangent} for quasi-static problems and
$\Keff$ of \eqref{eq:keff} for implicit dynamics.  With an exact tangent, an
exact linear solve, and $\alpha_k = 1$, convergence near the solution is
quadratic.  Each of those three assumptions is relaxed somewhere in Carina, and
each relaxation is the subject of a later chapter: approximate tangents in
Section~\ref{sec:tangent-strategies}, inexact linear solves in
Chapter~\ref{ch:inexact}, and $\alpha_k < 1$ immediately below.

\section{Globalization by line search}
\label{sec:linesearch}

Far from the solution the full Newton step can be worse than no step at all.
At large $\Delta t$ or a large load increment the predictor may sit far enough
from equilibrium that $\Uvec_k + \Delta\Uvec_k$ inverts elements, producing
$J \le 0$ and a non-finite residual.

Carina globalizes with Armijo backtracking on the merit function
\begin{equation}
  m(\alpha) = \tfrac12 \norm{\Rvec(\Uvec_k + \alpha \Delta\Uvec_k)}^2 ,
\end{equation}
accepting the first $\alpha \in \{1, \tau, \tau^2, \dots\}$ satisfying the
sufficient-decrease condition
\begin{equation}
  m(\alpha) \le \left(1 - 2 c\, \alpha\right) m(0),
  \label{eq:armijo}
\end{equation}
with backtrack factor $\tau = \tfrac12$ and $c = 10^{-4}$ by default.
Equation~\eqref{eq:armijo} is the standard Armijo condition specialized to this
merit function: since $\Delta\Uvec_k$ solves $\Kmat\Delta\Uvec = \Rvec$, the
directional derivative is $m'(0) = -\norm{\Rvec}^2 = -2m(0)$, and substituting
into $m(\alpha) \le m(0) + c\,\alpha\, m'(0)$ gives \eqref{eq:armijo} directly.

\begin{algorithm}[h]
\caption{Newton step with Armijo backtracking}
\label{alg:newton}
\begin{algorithmic}[1]
  \Require iterate $\Uvec_k$, residual $\Rvec_k$, tangent $\Kmat_k$,
           parameters $c$, $\tau$, budget $j_{\max}$
  \State $\eta_k \gets$ \Call{ForcingTerm}{$\cdot$}
    \Comment{Algorithm~\ref{alg:forcing}; $\eta_k = \varepsilon_{\mathrm{lin}}$ if disabled}
  \State solve $\Kmat_k \Delta\Uvec_k = \Rvec_k$ to relative tolerance $\eta_k$
  \State $m_0 \gets \tfrac12\norm{\Rvec_k}^2$,\quad $\alpha \gets 1$
  \For{$j = 1$ \textbf{to} $j_{\max}$}
    \State evaluate $\Rvec$ at $\Uvec_k + \alpha\Delta\Uvec_k$;\quad
           $m \gets \tfrac12\norm{\Rvec}^2$
    \If{$m \le (1 - 2c\alpha)\,m_0$}
      \State \Return $\Uvec_k + \alpha\Delta\Uvec_k$
        \Comment{Armijo satisfied}
    \EndIf
    \State $\alpha \gets \tau\alpha$
  \EndFor
  \State restore the pre-step material state
    \Comment{trials advanced it; see Remark~\ref{rem:restore}}
  \State \Return $\Uvec_k + \Delta\Uvec_k$
    \Comment{budget exhausted; outer failure machinery takes over}
\end{algorithmic}
\end{algorithm}

The line search is \emph{on by default}.  It costs roughly one extra residual
evaluation per iteration when the full step is already good, since $\alpha = 1$
is then accepted immediately, and it is the only thing standing between a bad
predictor and an inverted element.

\begin{remark}
\label{rem:restore}
If the backtracking budget is exhausted without satisfying
\eqref{eq:armijo}, Carina restores the pre-step material state and accepts the
full step, leaving the outer failure machinery of Section~\ref{sec:adaptive} to
handle it.  Restoring the state matters: the trial evaluations have already
advanced the internal variables of a path-dependent model, and accepting a step
on top of a state advanced by a rejected trial would silently corrupt the
history.
\end{remark}

\begin{remark}
An inexact Newton direction (Chapter~\ref{ch:inexact}) remains a descent
direction for $m$ as long as the linear residual is smaller than $\norm{\Rvec}$
--- that is, whenever the forcing term satisfies $\eta_k < 1$, which the clamp
of Section~\ref{sec:ew-safeguards} enforces.  The line search therefore
composes with inexact solves without modification.
\end{remark}

\section{Status tests}
\label{sec:status}

Termination is a \emph{tree} of composable tests rather than a fixed criterion,
following the NOX pattern.  Leaves test the residual ($\norm{\Rvec} <
\varepsilon_a$ or $\norm{\Rvec}/\norm{\Rvec_0} < \varepsilon_r$), the update
($\norm{\Delta\Uvec}$, absolute or relative to $\norm{\Uvec}$), the iteration
count, finiteness, divergence, or stagnation; interior nodes combine them with
\textsc{And}\,/\,\textsc{Or}.

Two properties of the combination rules are worth stating because algorithms
elsewhere depend on them.  In an \textsc{Or} node, \textsc{Converged} beats
\textsc{Failed} when both occur on the same iteration.  And a
\emph{finite-value} test is appended to every tree automatically, so a
\textsc{NaN} residual always terminates the solve regardless of what the user
wrote.

The tree is not merely a convenience.  Section~\ref{sec:ew-safeguards} needs
the residual norm the solve is actually aiming at, and reads it off this tree
--- taking the loosest child of an \textsc{Or} and the tightest of an
\textsc{And} --- rather than from a nominal tolerance that a user-supplied tree
may have overridden.

\chapter{Inexact Newton and forcing terms}
\label{ch:inexact}
\impl{src/linear\_solvers.jl}

\section{The over-solving problem}

Newton's step \eqref{eq:newton} is a linear system, and when it is solved by an
iterative method the accuracy of that solve is a free parameter.  An
\emph{inexact Newton method} accepts any $\Delta\Uvec_k$ satisfying
\begin{equation}
  \norm{\Rvec(\Uvec_k) - \Kmat(\Uvec_k)\Delta\Uvec_k} \le \eta_k \norm{\Rvec(\Uvec_k)},
  \label{eq:inexact}
\end{equation}
where $\eta_k \in [0,1)$ is the \emph{forcing term}.  Choosing $\eta_k$
constant and small --- the naive choice, and Carina's original behaviour ---
solves every linear system to the same tolerance regardless of how far the
nonlinear iteration still has to go.

That this wastes work is not an abstract concern.  On a $530{,}523$-DOF Newmark
step, conjugate gradients spent $181$, $205$, and $193$ iterations on the three
solves of a step while the nonlinear residual fell
\begin{equation*}
  1.61\times10^{4} \ \longrightarrow\ 1.34\times10^{2}
  \ \longrightarrow\ 7.79\times10^{-2} \ \longrightarrow\ 4.60\times10^{-8}.
\end{equation*}
The first two steps were computed to eight digits of linear accuracy and then
discarded by the next Newton correction.  The pattern on a quasi-static $J_2$
specimen is worse: roughly $350$ iterations on each of the first seven solves
while the residual crawls from $1.57\times10^3$ to $1.76\times10^2$, after
which the line search backtracks to $\alpha = \tfrac14$ and throws away most of
the step it just paid for.

\section{Eisenstat--Walker choice 2}

The remedy is to tie $\eta_k$ to the convergence actually being observed.
Carina implements choice~2 of Eisenstat and Walker~\cite{EisenstatWalker1996}:
\begin{equation}
  \eta_k^{\mathrm{EW}} = \gamma \left(\frac{\norm{\Rvec_k}}{\norm{\Rvec_{k-1}}}\right)^{\alpha},
  \qquad \gamma \in (0,1], \quad \alpha \in (1,2].
  \label{eq:ew}
\end{equation}
The reasoning is that the ratio of successive residual norms estimates how well
the linear model predicted the nonlinear decrease.  When the previous step
behaved as the model said it would, the ratio is small and a tighter solve is
warranted; when the model is doing badly, a tight solve is refining a direction
that is wrong anyway.  The default $\alpha = (1+\sqrt5)/2$ is the paper's, and
recovers superlinear convergence of the outer iteration.

On the first iteration of a step no ratio exists, and $\eta_0$ is taken from
the \code{initial} parameter.

\section{Three safeguards}
\label{sec:ew-safeguards}

Equation~\eqref{eq:ew} alone is not usable.  Carina composes it with three
guards, applied in order.

\paragraph{1. Against $\eta$ collapsing too fast.}  A single unusually good
step can drive \eqref{eq:ew} to a very small value that the local convergence
rate does not justify, wasting on the next solve exactly what the method exists
to save.  Eisenstat and Walker's own safeguard raises $\eta$ back toward its
predecessor whenever that predecessor was itself large:
\begin{equation}
  \eta_k \leftarrow \max\left(\eta_k^{\mathrm{EW}},\ \gamma\,\eta_{k-1}^{\alpha}\right)
  \qquad\text{whenever}\qquad \gamma\,\eta_{k-1}^{\alpha} > 0.1 .
  \label{eq:ew-guard}
\end{equation}
The threshold $0.1$ is the paper's and is not a tuning parameter in Carina.  It
nonetheless has a consequence for the parameter that \emph{is} tunable, which
Section~\ref{sec:etamax} takes up.

\paragraph{2. Against over-solving the final step.}  Near convergence the ratio
in \eqref{eq:ew} becomes tiny and demands far more linear accuracy than the
nonlinear stopping test will ever reward.  Kelley's guard~\cite{Kelley1995}
bounds $\eta$ from below by what the target actually requires:
\begin{equation}
  \eta_k \leftarrow \max\left(\eta_k,\ \frac{s\,\tau}{\norm{\Rvec_k}}\right),
  \label{eq:kelley}
\end{equation}
where $\tau$ is the residual norm the nonlinear test is aiming at and $s$ a
safety factor, $\tfrac12$ by default.  Carina reads $\tau$ from the status test
tree of Section~\ref{sec:status}, so a deck's own \code{termination} block
governs rather than a nominal tolerance it may have overridden.

\paragraph{3. A hard clamp.}  Finally
\begin{equation}
  \eta_k \leftarrow \min\left(\max\left(\eta_k,\ \varepsilon_{\mathrm{lin}}\right),\ \eta_{\max}\right),
  \label{eq:clamp}
\end{equation}
where $\varepsilon_{\mathrm{lin}}$ is the deck's own linear tolerance.

The order matters, and Algorithm~\ref{alg:forcing} fixes it: the two lower
bounds are applied before the clamp, so neither can push $\eta$ outside
$[\varepsilon_{\mathrm{lin}}, \eta_{\max}]$.

\begin{algorithm}[h]
\caption{Forcing term at Newton iteration $k$}
\label{alg:forcing}
\begin{algorithmic}[1]
  \Require $\norm{\Rvec_k}$; previous $\norm{\Rvec_{k-1}}$ and $\eta_{k-1}$
           if any; target $\tau$; parameters $\gamma, \alpha, \eta_{\max},
           \eta_0, s$; linear tolerance $\varepsilon_{\mathrm{lin}}$
  \If{no previous residual}
    \State $\eta \gets \eta_0$
      \Comment{first iteration of a step}
  \Else
    \State $\eta \gets \gamma\left(\norm{\Rvec_k}/\norm{\Rvec_{k-1}}\right)^{\alpha}$
    \If{$\gamma\,\eta_{k-1}^{\alpha} > 0.1$}
      \State $\eta \gets \max\left(\eta,\ \gamma\,\eta_{k-1}^{\alpha}\right)$
        \Comment{safeguard 1: anti-collapse}
    \EndIf
  \EndIf
  \If{$s > 0$ \textbf{and} $\norm{\Rvec_k} > 0$}
    \State $\eta \gets \max\left(\eta,\ s\,\tau/\norm{\Rvec_k}\right)$
      \Comment{safeguard 2: over-solve guard}
  \EndIf
  \State $\eta_k \gets \min\left(\max\left(\eta,\ \varepsilon_{\mathrm{lin}}\right),\ \eta_{\max}\right)$
    \Comment{safeguard 3: clamp}
  \State \Return $\eta_k$
\end{algorithmic}
\end{algorithm}

\begin{proposition}[Safety]
\label{prop:safety}
Under \eqref{eq:clamp}, a forcing term can only loosen a linear solve relative
to the fixed-tolerance baseline, never tighten one.
\end{proposition}

\noindent This is the property that makes the method switchable on without
re-validating a model.  The lower clamp guarantees the final Newton iterations
run at exactly the tolerance the user requested, so the converged solution is
as accurate as it was before and an A/B comparison against a fixed tolerance
compares like with like.  Measurements bear this out: across CPU and A100 runs
in both implicit regimes, $\norm{\Uvec}_\infty$ agrees with the
fixed-tolerance baseline to every printed digit at every output stop.

\begin{remark}
Preconditioned conjugate gradients stops on the preconditioned residual
$\Prec$-norm relative to its own initial value, not on the unpreconditioned
ratio of \eqref{eq:inexact} for which the theory is stated.  The forcing term
is a heuristic in either norm; Proposition~\ref{prop:safety} is what bounds the
damage when the two disagree.
\end{remark}

\section{Choosing \texorpdfstring{$\eta_{\max}$}{eta max}}
\label{sec:etamax}

The literature's default is $\eta_{\max} = 0.9$.  Carina's is $0.2$, and the
reason is structural rather than empirical fitting.

Measurements across a CPU quasi-static sweep and two A100 regimes show that the
\emph{total} linear work is nearly invariant in $\eta_{\max}$: $63.1$k, $62.9$k
and $63.6$k CG iterations at $\eta_{\max} = 0.1$, $0.2$, $0.5$ respectively on
the $J_2$ specimen, and $1019$ versus $1011$ on the A100 Newmark problem.  What
$\eta_{\max}$ actually controls is how many \emph{Newton} iterations are spent
buying that work --- $413$ versus $547$ on the same CPU sweep.

The mechanism is safeguard~\eqref{eq:ew-guard}, which activates when
$\gamma\eta^{\alpha} > 0.1$.  At the default $\alpha$:
\begin{center}
\begin{tabular}{@{}lll@{}}
\toprule
$\eta_{\max}$ & $\eta_{\max}^{\alpha}$ & Behaviour \\
\midrule
$0.2$ & $0.076$ & below the knee; the residual ratio drives $\eta$ \\
$0.5$ & $0.326$ & above it; the safeguard pins $\eta$ high for several iterations \\
$0.9$ & $0.842$ & far above \\
\bottomrule
\end{tabular}
\end{center}
So $0.2$ is the largest round value that keeps Eisenstat--Walker's
anti-collapse safeguard out of the way of Eisenstat--Walker's own adaptivity.
Two different problems on two different architectures selected it
independently.

\section{What the method buys}

A forcing term trades linear work for nonlinear work.  It therefore pays in
proportion to how much of a step is the linear solve --- most on the
matrix-free GPU paths of Chapter~\ref{ch:matrixfree}, where the step
essentially \emph{is} the stiffness matrix--vector product, and least on
assembled CPU runs where residual assembly and the line search already
dominate.  Measured with the defaults:

\begin{center}
\begin{tabular}{@{}lcc@{}}
\toprule
Configuration & CG iterations & Per-step wall \\
\midrule
CPU quasi-static, $J_2$, $57$k DOF        & $2.47\times$ fewer & $1.30\times$ \\
A100 Newmark, $530$k DOF, CG + Jacobi     & $2.27\times$ fewer & $1.45\times$ \\
A100 quasi-static, $530$k DOF, CG + AMG   & $2.05\times$ fewer & $1.59\times$ \\
\bottomrule
\end{tabular}
\end{center}

The quasi-static figure is on top of the roughly twofold advantage algebraic
multigrid already holds over Jacobi on that problem.

\section{Interaction with lagged preconditioners}

The hierarchy staleness detector of Section~\ref{sec:amg-lagging} compares
conjugate-gradient iteration counts across solves.  Under a forcing term those
counts are taken at \emph{different} tolerances and are not comparable as raw
numbers.

Simply ignoring loosened solves does not work, and the reason is
guard~\eqref{eq:kelley}: it can raise $\eta$ enough that a run converges
without ever solving at the deck's tolerance, so a detector that only counted
full-tolerance solves would never fire at all.  Since that detector is the only
rebuild trigger on the quasi-static path, the hierarchy would be built once at
the reference configuration and kept for an entire load history --- precisely
the failure Section~\ref{sec:nullspace} exists to prevent.

Carina instead rescales.  Because a Krylov residual falls geometrically, the
iteration count is roughly proportional to the number of digits of reduction
requested, so
\begin{equation}
  \hat{m}_k = m_k \cdot \frac{\log \varepsilon_{\mathrm{lin}}}{\log \eta_k}
  \label{eq:rescale}
\end{equation}
converts a count $m_k$ taken at tolerance $\eta_k$ into the count the same
solve would have cost at the deck's tolerance.  At $\eta_k =
\varepsilon_{\mathrm{lin}}$ this is exactly the identity, so a run with no
forcing term feeds the detector precisely the numbers it always did.  The
estimate is coarse, but the detector's threshold --- a factor of three ---
is coarser.

\chapter{Quasi-Newton and gradient methods}
\label{ch:quasinewton}
\impl{src/nonlinear\_solvers.jl, src/linear\_solvers.jl}

Three alternatives to Newton are implemented, all matrix-free in the sense that
none forms or factors a tangent.  They exist for cases where assembling
$\Kmat$ is unattractive, and each has a regime where it is the right choice
and a regime where it fails.

\section{L-BFGS}
\label{sec:lbfgs}

Limited-memory BFGS builds an implicit approximation to $\Kmat^{-1}$ from the
last $m$ secant pairs
\begin{equation}
  \vect{s}_k = \Uvec_{k+1} - \Uvec_k,
  \qquad
  \vect{y}_k = \Rvec_k - \Rvec_{k+1},
  \qquad
  \rho_k = \frac{1}{\vect{y}_k^{\mathsf{T}}\vect{s}_k},
\end{equation}
and applies it through the two-loop recursion: given $\vect{q} = \Rvec_k$,
sweep newest to oldest computing $a_i = \rho_i \vect{s}_i^{\mathsf{T}}\vect{q}$
and $\vect{q} \leftarrow \vect{q} - a_i \vect{y}_i$; apply an initial inverse
Hessian $\tensor{H}_0$; then sweep oldest to newest correcting with
$b_i = \rho_i \vect{y}_i^{\mathsf{T}}\vect{d}$ and $\vect{d} \leftarrow
\vect{d} + (a_i - b_i)\vect{s}_i$.  The cost is $O(mn)$ per iteration with no
matrix at all.

The choice of $\tensor{H}_0$ matters more here than the textbook presentation
suggests.  Carina uses Barzilai--Borwein scaling from the most recent pair,
\begin{equation}
  \tensor{H}_0 = \gamma_0 \Ident,
  \qquad
  \gamma_0 = \frac{\vect{s}_k^{\mathsf{T}}\vect{y}_k}{\vect{y}_k^{\mathsf{T}}\vect{y}_k},
\end{equation}
once history exists, and a Jacobi diagonal before then.  The Jacobi fallback is
not cosmetic: on the first step of an implicit dynamic problem there is no
history, and $c_M = 1/(\beta\Delta t^2)$ can be enormous
(equation~\ref{eq:cM}), so $\tensor{H}_0 = \Ident$ would produce a first step
wrong by many orders of magnitude.

\paragraph{Where it works, and where it does not.}  L-BFGS is the fastest
option Carina has for implicit \emph{dynamics} on GPU, and it fails outright on
quasi-static problems --- stalling roughly seven orders short of tolerance, on
CPU and GPU alike, at any history size.

Section~\ref{sec:newmark} explains why.  At small $\Delta t$ the effective
operator $\Keff = \Kmat + c_M\Mmat$ is mass-dominated and strongly diagonally
dominant, which is exactly the structure a rank-$m$ update models well.
Quasi-static has no mass term: the operator is the bare stiffness, whose
spectrum spans the mesh-refinement range, and a rank-$10$ model cannot
represent it.  This is a statement about the operator, not about the
implementation or the device --- which is why no amount of history size
rescues it.

\section{Nonlinear conjugate gradient}

Preconditioned Polak--Ribi\`ere$^+$.  With $\vect{g}_k = \Prec\Rvec_k$ the
preconditioned gradient,
\begin{equation}
  \beta_k = \max\left(0,\
  \frac{\Rvec_k^{\mathsf{T}}\vect{g}_k - \Rvec_k^{\mathsf{T}}\vect{g}_{k-1}}
       {\Rvec_{k-1}^{\mathsf{T}}\vect{g}_{k-1}}\right),
  \qquad
  \vect{d}_k = \vect{g}_k + \beta_k \vect{d}_{k-1}.
\end{equation}
The clamp at zero is the ``$+$'' variant and amounts to restarting with a
steepest-descent direction whenever the formula would turn negative, which
guarantees $\vect{d}_k$ remains a descent direction.  Carina restarts
additionally when successive gradients lose orthogonality beyond a tolerance,
and optionally at a fixed interval.

Only Jacobi preconditioning is available at this level; asking for anything
stronger is rejected rather than silently downgraded, since a preconditioner
that silently did nothing would present as a mysterious convergence problem.

\section{Preconditioned steepest descent}

The direction is $\vect{d}_k = \Prec\Rvec_k$, with an Armijo line search on the
total potential energy rather than on the residual norm:
\begin{equation}
  \Pi(\Uvec + \alpha\vect{d}) \le \Pi(\Uvec) + c\,\alpha\, \nabla\Pi\cdot\vect{d},
  \qquad \nabla\Pi\cdot\vect{d} = -\Rvec^{\mathsf{T}}\vect{d}.
\end{equation}
Using the energy is available here because the directional derivative is exact
and cheap, and $\Pi$ is the quantity the equilibrium problem actually
minimizes.  This is the slowest and most robust option: linear convergence, but
it makes progress on problems where the others stall.

\chapter{Krylov methods}
\label{ch:krylov}
\impl{src/linear\_solvers.jl}

\section{Why conjugate gradients}

Every operator Carina hands to a linear solver is symmetric positive definite:
the tangent stiffness of a hyperelastic material is the Hessian of a stored
energy, and the Newmark effective operator \eqref{eq:keff} adds $c_M\Mmat$ with
$c_M > 0$, which can only improve matters.  Conjugate gradients is therefore
the method of choice, and the input parser maps every iterative alias
--- \code{cg}, \code{krylov}, \code{minres} --- onto it.

The assumption is checked rather than trusted.  Carina measures
$\norm{\Kmat - \Kmat^{\mathsf{T}}}_\infty / \norm{\Kmat}_\infty$ once per run
and reports if the tangent is not symmetric to roundoff, since CG applied to a
non-symmetric operator does not merely converge slowly --- it converges to the
wrong answer, silently.  The measured value on the assembled path is
$\sim 10^{-16}$.

\section{Preconditioned conjugate gradients}
\label{sec:pcg}

With a symmetric positive definite preconditioner $\Prec \approx \Kmat^{-1}$,
the iteration is
\begin{align}
  \vect{z}_i &= \Prec \vect{r}_i, \qquad \rho_i = \vect{r}_i^{\mathsf{T}}\vect{z}_i, \\
  \vect{p}_i &= \vect{z}_i + (\rho_i/\rho_{i-1})\, \vect{p}_{i-1}, \\
  \alpha_i   &= \rho_i / \left(\vect{p}_i^{\mathsf{T}}\Kmat\vect{p}_i\right), \\
  \Uvec_{i+1} &= \Uvec_i + \alpha_i \vect{p}_i,
  \qquad
  \vect{r}_{i+1} = \vect{r}_i - \alpha_i \Kmat \vect{p}_i .
\end{align}
One operator application and one preconditioner application per iteration.

The error bound
\begin{equation}
  \norm{\vect{e}_i}_{\Kmat} \le 2\left(\frac{\sqrt{\kappa}-1}{\sqrt{\kappa}+1}\right)^{i}
                                 \norm{\vect{e}_0}_{\Kmat},
  \qquad \kappa = \kappa(\Prec\Kmat),
  \label{eq:cg-bound}
\end{equation}
is what every preconditioner in Chapter~\ref{ch:preconditioners} exists to
improve, and it is also the justification for the rescaling
\eqref{eq:rescale}: \eqref{eq:cg-bound} is geometric in $i$, so iterations are
proportional to digits of reduction requested.

\paragraph{Stopping.}  Carina terminates on the preconditioned residual
$\Prec$-norm relative to its own initial value, $\sqrt{\rho_i} \le \eta
\sqrt{\rho_0}$, because $\rho_i = \vect{r}_i^{\mathsf{T}}\Prec\vect{r}_i$ is
already computed each iteration and needs no extra inner product.  The
consequence for the forcing-term theory is noted in
Chapter~\ref{ch:inexact}.

\paragraph{Cold start.}  The workspace solution is zeroed before every solve
rather than warm-started from the previous Newton iteration.  A stale
warm-start is not merely unhelpful here: the operator has changed between
Newton iterations, so the previous solution is the answer to a different
question, and beginning from it can leave the residual larger than beginning
from zero.

\section{Device-resident conjugate gradients}
\label{sec:device-cg}

On GPU the entire iteration runs on the device.  This is not a detail of
implementation but a change in what dominates the cost: a host-resident CG must
transfer $\rho_i$ and $\alpha_i$ back to the host each iteration to make the
scalar decisions, and each transfer is a synchronization point that idles the
device for longer than the arithmetic it guards.

The dot products are reduced on-device by a two-stage tree --- a fixed number
of workgroups each reducing a strided slice, then a single workgroup reducing
the partials --- and the scalars stay in device memory.  The only quantity that
must reach the host is the convergence flag.

\section{Direct solvers}
\label{sec:gpu-cholesky}

A sparse $LDL^{\mathsf{T}}$ factorization remains the right choice for small to
medium problems on CPU: it needs no tuning and cannot fail to converge.  It is
unavailable on GPU, and asking for it there is a hard error rather than a
silent fallback.

One hybrid exists.  For a linear-elastic model the stiffness is constant, so
Carina factors $\Kmat$ once on the host, uploads the Cholesky factor to the
device, and thereafter solves by two device triangular solves.  This gives
direct-solver robustness with no host transfer per step, and applies only
because linearity makes the factorization reusable for the entire run.

\chapter{Preconditioners}
\label{ch:preconditioners}
\impl{src/linear\_solvers.jl}

A preconditioner is an approximation $\Prec \approx \Kmat^{-1}$ that is cheap
to apply and reduces $\kappa(\Prec\Kmat)$ in \eqref{eq:cg-bound}.  It must be
symmetric positive definite for conjugate gradients to remain valid.  It need
not be accurate: Remark~\ref{ch:constitutive} applies here too --- the
preconditioner shapes the path, never the fixed point.

\section{Jacobi}

$\Prec = \Dmat^{-1}$ with $\Dmat = \diag(\Kmat)$.  One vector scaling per
application, trivially parallel, always available.  It removes the
conditioning that comes from heterogeneous element sizes and material
properties, and nothing else --- in particular it does nothing about the
$(L/h)^2$ growth with mesh refinement that dominates quasi-static problems.

Its unusual effectiveness in implicit dynamics at small $\Delta t$ is a direct
consequence of \eqref{eq:cM}: when $c_M\Mmat$ dominates $\Keff$, and $\Mmat$ is
nearly diagonal for the element types in use, the diagonal captures most of the
operator.

\section{Incomplete \texorpdfstring{$LDL^{\mathsf{T}}$}{LDLT}}

An incomplete factorization $\Kmat \approx \tensor{L}\Dmat\tensor{L}^{\mathsf{T}}$
that drops fill outside a chosen sparsity pattern, applied by two triangular
solves.  Substantially stronger than Jacobi for ill-conditioned systems, and
the natural choice for $J_2$ plasticity on non-uniform meshes.

It is CPU-only, and the obstruction is structural rather than a missing port: a
triangular solve is inherently sequential along each dependency chain, which is
the one access pattern a GPU cannot exploit.  This is why the GPU path needs
either a polynomial preconditioner or multigrid, both of which reduce to
matrix--vector products.

\section{Chebyshev polynomial preconditioner}
\label{sec:chebyshev}

A polynomial preconditioner approximates $\Kmat^{-1}$ by a polynomial in
$\Kmat$, so applying it costs only matrix--vector products --- which makes it
GPU-friendly where triangular solves are not.

Carina implements Chebyshev polynomials of the \emph{fourth kind} with optimal
weights, following Lottes~\cite{Lottes2022} as implemented in Ifpack2.  This
variant has three properties that matter here:

\begin{itemize}
  \item It requires only $\lambda_{\max}$, not the eigenvalue \emph{ratio}.
        Standard Chebyshev needs a lower bound $\lambda_{\min}$ as well, and a
        poor estimate of it degrades the polynomial badly; $\lambda_{\min}$ is
        also the harder of the two to estimate.
  \item Jacobi scaling is baked into the recurrence, so $\Dmat^{-1}$ is applied
        at each stage rather than as a separate wrapper.
  \item The optimal weights make the resulting operator symmetric positive
        definite \emph{by construction}, so it is admissible as a CG
        preconditioner without the squaring trick --- costing $k$ matrix--vector
        products per application instead of $2k$.
\end{itemize}

With $\sigma = 1/\lambda_{\max}$ and weights $\beta_i$ tabulated for degrees
$1$ through $16$:

\begin{algorithm}[h]
\caption{Fourth-kind Chebyshev preconditioner application, degree $k$}
\label{alg:chebyshev}
\begin{algorithmic}[1]
  \Require $\vect{b}$, operator $\Kmat$, $\Dmat^{-1}$, $\sigma = 1/\lambda_{\max}$,
           weights $\beta_1,\dots,\beta_k$
  \Ensure $\vect{y} \approx \Kmat^{-1}\vect{b}$
  \State $\vect{z} \gets \tfrac43\sigma\,\Dmat^{-1}\vect{b}$,\quad
         $\vect{x} \gets \vect{z}$,\quad
         $\vect{y} \gets \beta_1\vect{z}$
  \For{$i = 1$ \textbf{to} $k-1$}
    \State $\gamma_i \gets \dfrac{2i-1}{2i+3}$,\qquad
           $\rho_i \gets \dfrac{8i+4}{2i+3}\,\sigma$
    \State $\vect{z} \gets \rho_i\,\Dmat^{-1}\left(\vect{b} - \Kmat\vect{x}\right)
                            + \gamma_i\vect{z}$
      \Comment{one operator application}
    \State $\vect{x} \gets \vect{x} + \vect{z}$
      \Comment{raw fourth-kind iterate}
    \State $\vect{y} \gets \vect{y} + \beta_{i+1}\vect{z}$
      \Comment{weighted accumulation}
  \EndFor
  \State \Return $\vect{y}$
\end{algorithmic}
\end{algorithm}

\noindent Note that $\vect{x}$ and $\vect{y}$ are distinct: $\vect{x}$ carries
the unweighted recurrence that generates the polynomial basis, while $\vect{y}$
accumulates the optimally weighted combination that is actually returned.
Collapsing them --- returning $\vect{x}$ --- gives the ordinary fourth-kind
polynomial without the optimal weights, and loses the positive-definiteness
that makes the operator admissible to conjugate gradients.

\paragraph{Spectral bound.}  $\lambda_{\max}$ of $\Dmat^{-1}\Kmat$ is estimated
by ten power iterations followed by a Rayleigh quotient, then inflated by a
safety factor of $1.1$.  The inflation is deliberate and one-sided:
overestimating $\lambda_{\max}$ makes the polynomial conservative and merely
slower, whereas underestimating it puts part of the spectrum outside the
interval the polynomial damps, and the preconditioner amplifies those modes
instead of removing them.

\paragraph{Symmetric scaling.}  On the assembled path the polynomial is built
on $\Dmat^{-1/2}\Kmat\Dmat^{-1/2}$ rather than on $\Kmat$, for the reason given
in Section~\ref{sec:elimination}: penalty-scale entries give $\Kmat$
eigenvalues of order $10^{15}$ while the symmetrically scaled operator has
eigenvalues of order one, and a polynomial fitted to the former is useless.
The preconditioner is then $\Prec = \Dmat^{-1/2} p_k(\tensor{S})\, \Dmat^{-1/2}$,
which is symmetric positive definite whenever $p_k$ is.

\chapter{Smoothed-aggregation multigrid}
\label{ch:amg}
\impl{src/gpu\_amg.jl, src/linear\_solvers.jl}

\section{Why multigrid}

Every preconditioner of Chapter~\ref{ch:preconditioners} is local: it couples
each unknown to its immediate algebraic neighbours.  Local operators damp
high-frequency error efficiently and low-frequency error not at all, which is
exactly why $\kappa(\Kmat)$ grows like $(L/h)^2$ under refinement --- the
smoothest error modes are the ones a local operator cannot see.

Multigrid removes what smoothing cannot by representing the smooth error on a
coarser problem, where it is no longer smooth relative to the grid.  The payoff
is an iteration count that is nearly \emph{independent} of problem size.  On
a $530$k-DOF torsion problem, algebraic multigrid holds $5$--$17$ conjugate
gradient iterations across a range of $\Delta t$ over which Jacobi grows from
$30$ to $442$.

Carina uses \emph{algebraic} multigrid: the hierarchy is built from the matrix
alone, with no geometric coarse mesh, which is what makes it applicable to
unstructured meshes.

\section{The near-nullspace}
\label{sec:nullspace}

Algebraic multigrid needs to know which error modes the smoother will not
remove, so that the coarse space can represent them.  For elasticity these are
the modes with no strain energy: the six rigid-body modes --- three
translations and three rotations.  They form the columns of the
\emph{near-nullspace} $\Bnull \in \mathbb{R}^{3n_{\mathrm{node}} \times 6}$,
built from nodal coordinates about the centroid $\bar{\vect{x}}$:
\begin{equation}
  \Bnull = \left[\
  \begin{array}{ccc|ccc}
    1 & 0 & 0 & 0     & \hat{z}  & -\hat{y} \\
    0 & 1 & 0 & -\hat{z} & 0     & \hat{x}  \\
    0 & 0 & 1 & \hat{y}  & -\hat{x} & 0
  \end{array}\ \right],
  \qquad \hat{\vect{x}} = \vect{x} - \bar{\vect{x}},
\end{equation}
one such block per node, restricted to the free degrees of freedom.

\paragraph{Current configuration, not reference.}  Carina builds $\Bnull$ from
the \emph{current} nodal coordinates $\vect{X} + \disp$ at every hierarchy
build, not from the reference configuration.  At finite deformation the true
near-nullspace of $\Kmat(\Uvec)$ consists of rigid-body motions about the
\emph{deformed} body; a nullspace frozen at the reference configuration
describes rotations the body no longer has.  The consequence is not subtle:
on a twisted-bar run, a frozen reference nullspace degraded the iteration count
from $63$ to over $400$ as the bar rotated, while the hierarchy build count
stayed at one and nothing in the log indicated a problem.

\section{Hierarchy construction}

Each level is built in four stages.

\paragraph{1. Strength of connection.}  A symmetric strength measure filters
$\Kmat$ into a graph $\tensor{S}$ retaining only algebraically strong
couplings, so that aggregation follows the directions in which error actually
propagates.

\paragraph{2. Aggregation.}  Standard aggregation partitions the nodes of
$\tensor{S}$ into disjoint aggregates, each becoming one coarse node.

\paragraph{3. Tentative prolongator.}  For each aggregate $g$, the restriction
of the near-nullspace to that aggregate's rows is factored,
\begin{equation}
  \Bnull|_g = \tensor{Q}_g \tensor{R}_g,
\end{equation}
by a thin QR.  The orthonormal $\tensor{Q}_g$ becomes the aggregate's block of
the tentative prolongator $\tensor{T}$, and $\tensor{R}_g$ becomes the coarse
near-nullspace $\Bnull_c$ for the next level.  This construction guarantees
\begin{equation}
  \Bnull = \tensor{T}\Bnull_c ,
\end{equation}
i.e.\ the coarse space represents every near-null mode \emph{exactly} --- which
is the whole point of supplying $\Bnull$.

\begin{remark}
An aggregate with fewer nodes than $\Bnull$ has columns is rank deficient, and
its QR yields fewer than six usable columns.  The implementation handles this
by giving such an aggregate $r = \min(\abs{g}, 6)$ columns and leaving the
remainder structurally empty, rather than producing near-zero columns that
would inflate every downstream sparse product.
\end{remark}

\paragraph{4. Prolongator smoothing.}  The tentative prolongator is improved by
one damped Jacobi sweep,
\begin{equation}
  \Pmat = \left(\Ident - \frac{4}{3\lambda_{\max}}\Dmat^{-1}\Kmat\right)\tensor{T},
\end{equation}
which lowers the energy of the coarse basis functions and is what distinguishes
\emph{smoothed} aggregation from plain aggregation.  Restriction is
$\Rrest = \Pmat^{\mathsf{T}}$, and the coarse operator is the Galerkin triple
product
\begin{equation}
  \Kmat_c = \Rrest\, \Kmat\, \Pmat .
  \label{eq:galerkin}
\end{equation}

\begin{remark}[Memory]
Evaluated naively, \eqref{eq:galerkin} preallocates from an up-front
nonzero-count estimate that is wildly pessimistic for a six-column
near-nullspace --- demanding tens of gigabytes at $1.5$M DOF for a product
whose true size is about $4$~GB.  Carina evaluates the fine-level product in
column slabs of $\Pmat$, so peak transient memory is one slab's product rather
than the whole estimate, and hands the much smaller coarse problem to the
standard routine for the remaining levels.
\end{remark}

\section{The V-cycle}

One V($\nu,\nu$)-cycle applied to a residual $\vect{r}$ produces a correction
$\vect{z}$:

\begin{algorithm}[h]
\caption{V-cycle, level $\ell$}
\begin{algorithmic}[1]
  \STATE $\vect{x} \gets \vect{0}$
  \STATE pre-smooth $\nu$ times on $\Kmat^{\ell}\vect{x} = \vect{b}^{\ell}$
  \STATE $\vect{b}^{\ell+1} \gets \Rrest^{\ell}\left(\vect{b}^{\ell} - \Kmat^{\ell}\vect{x}\right)$
  \STATE recurse to level $\ell+1$, or solve directly if coarsest
  \STATE $\vect{x} \gets \vect{x} + \Pmat^{\ell}\vect{x}^{\ell+1}$
  \STATE post-smooth $\nu$ times
\end{algorithmic}
\end{algorithm}

The smoother is damped Jacobi with $\omega = 4/(3\lambda_{\max})$, the same
weight used for prolongator smoothing, with $\lambda_{\max}$ estimated per
level by power iteration.  The coarsest level is solved by a dense
pseudo-inverse, formed once at build time; at that size a dense solve is
cheaper than another level of hierarchy.

Symmetry of the cycle --- equal pre- and post-smoothing, and
$\Rrest = \Pmat^{\mathsf{T}}$ --- is what makes the resulting operator
symmetric, hence admissible as a conjugate-gradient preconditioner.

\section{GPU execution}

The hierarchy is built on the host, where the sparse pattern already lives, and
converted to device CSR.  The cycle then runs entirely on the device.  The fine
level is special: rather than storing $\Kmat$ on the device, its
matrix--vector products go through the matrix-free action of
Chapter~\ref{ch:matrixfree}, so the fine matrix is never formed there at all.
Only the coarse levels --- orders of magnitude smaller --- are stored
explicitly.

\section{Lagging and staleness}
\label{sec:amg-lagging}

Hierarchy setup costs seconds at $500$k DOF, far more than a single V-cycle
application, so the hierarchy is \emph{lagged}: built once and reused across
Newton iterations and time steps.  Reuse is safe because the hierarchy is a
preconditioner --- an out-of-date one costs iterations, never accuracy.

Two triggers force a rebuild.  A change in $c_M$ means the operator itself has
changed (a new $\Delta t$).  And a staleness detector watches the
conjugate-gradient iteration count: it latches a baseline on the first solve
after a build and requests a rebuild when a later count exceeds three times
that baseline, subject to a floor so that small counts do not thrash.

On the quasi-static path this detector is the \emph{only} rebuild trigger,
because $c_M = 0$ there and the first test never fires.  That is what makes the
interaction with forcing terms consequential, and it is treated in
Section~\ref{sec:etamax}'s neighbouring discussion at the end of
Chapter~\ref{ch:inexact}.

\section{When multigrid is the right tool}
\label{sec:amg-when}

AMG targets operators whose difficulty is the smooth end of the spectrum:
quasi-static problems and moderate-$\Delta t$ dynamics.  It is not universally
better, and two limits are worth naming.

At small $\Delta t$ in implicit dynamics, $\Keff$ is mass-dominated
(Section~\ref{sec:newmark}) and already well conditioned.  There is little
smooth error for a coarse grid to remove, and measurements on two GPU
architectures confirm the V-cycle does not repay its per-application cost:
Jacobi is faster.

At very large $\Delta t$ on violently dynamic problems, the Newmark predictor
can overshoot into near-inverted configurations whose tangent is indefinite.
Conjugate gradients is then invalid regardless of preconditioner, and explicit
integration is the appropriate method.

\chapter{Matrix-free operators}
\label{ch:matrixfree}
\impl{src/physics.jl, src/two\_phase\_action.jl}

\section{The action, not the matrix}

Conjugate gradients never needs $\Kmat$ --- only its action $\Kmat\vect{v}$ on
a sequence of vectors.  A matrix-free operator supplies that action by
assembling directly:
\begin{equation}
  \left(\Kmat\vect{v}\right)_a
  = \int_{\Omega_0} \Grad N_a \cdot \fourth{A} \cdot \Grad N_b \, \mathrm{d}V \ v_b
  = \int_{\Omega_0} \Grad N_a \dbldot \delta\Piola(\Grad \vect{v}) \, \mathrm{d}V,
  \label{eq:action}
\end{equation}
where $\delta\Piola$ is the directional derivative of the stress in the
direction $\Grad\vect{v}$ --- equation~\eqref{eq:nh-action} for neo-Hookean.
The right-hand form is the one that matters: the tangent's action on
\emph{one} direction never requires the $81$-component $\fourth{A}$ to be
formed, so the closed-form derivative costs a few $3\times3$ products per
quadrature point.

Three things follow.  Memory drops from $O(\text{nnz})$ to $O(n)$ --- measured
at $5.83$~GB down to $0.245$~GB on a $530$k-DOF problem, which is the
difference between fitting on a GPU and not.  There is no assembly phase and no
sparse pattern to build.  And the operator is exactly as current as the
displacement it is evaluated at, so no staleness question arises.

The cost is that each application redoes the constitutive work.  Whether that
is a good trade depends on the arithmetic intensity of the element kernel
against the bandwidth of storing and re-reading a sparse matrix; on GPU it is
decisively favourable.

\section{Operator fusion}

The Newmark system needs $\Keff\vect{v} = \Kmat\vect{v} + c_M\Mmat\vect{v}$.
Computed as two separate assemblies, the element loop runs twice and each pass
gathers the same nodal data.  Fusing them into a single kernel that accumulates
both contributions per quadrature point halves the gather traffic, and the
gather is what the operator is bound by.

The same reasoning produces a diagonal-only kernel for the Jacobi
preconditioner: $\diag(\Keff)$ is assembled directly rather than extracted from
an operator that is never formed.

\section{Where the time goes}

Profiling the fused Newmark action on an A100 at $530$k DOF attributes $67\%$
of device time to the stiffness matrix--vector product, with the device busy
$72\%$ of the step.  The step, to a good approximation, \emph{is} the matvec.

That single fact organizes the optimization work.  It is why atomic contention
is not the bottleneck --- ablating the atomics changes nothing measurable ---
and why the loss ordering across architectures tracks effective cache size
rather than arithmetic throughput.  It is also why the forcing terms of
Chapter~\ref{ch:inexact} pay so much better on the GPU paths than on assembled
CPU runs: on this path, removing a conjugate-gradient iteration removes almost
a whole iteration's worth of wall time.

\section{Two-phase assembly}

The default kernel assigns one thread per element: gather the element's nodal
values, loop the quadrature points serially, scatter the results with atomic
adds.  The working set held across that quadrature loop is what limits
occupancy.

The two-phase form splits it.  Phase one assigns one thread per
(element, quadrature point) pair, writing each contribution to a disjoint slot
of a staging array --- eight times the threads, an eighth of the per-thread
state, and no atomics.  Phase two assigns one thread per node and sums the
staging slots incident on it through an inverse adjacency built once on the
host, again without atomics.

The transformation is sound and it is implemented, but measurement did not
support it: $0.66$ to $0.80$ times the throughput of the one-thread-per-element
kernel on all three architectures tested.  The extra staging traffic exceeds
what the occupancy gain returns.  It is documented here because the reasoning
is correct and the conclusion is specific to the current data layout --- an
element-blocked field layout could change it.

\section{Mixed precision}
\label{sec:mixed-precision}

The V-cycle of Chapter~\ref{ch:amg} applies the matrix-free action five times
per call at $\nu = 2$, which on a consumer GPU --- where double precision runs
at $1/32$ rate --- dominates the iteration.  Carina evaluates the fine-level
\emph{smoother} action in single precision while the operator conjugate
gradients solves against stays in double.

This is legitimate, and the argument is the one from
Remark~\ref{ch:constitutive}.  A preconditioner may be any symmetric positive
definite approximation; reducing its precision changes which approximation it
is, not what the iteration converges to.  What would \emph{not} be legitimate
is reducing the precision of the operator itself, which would change the fixed
point.  The reduced-precision V-cycle also remains a fixed linear operator
across applications, so ordinary conjugate gradients --- rather than a flexible
variant --- stays valid.

The gate is a property of the model, not a user request: Carina enables the
reduced-precision path only for models that demonstrably narrow to single
precision, and verifies this rather than assuming it.  A model whose action
silently falls back to double would otherwise present as a preconditioner that
was never actually accelerated.


\backmatter
\bibliographystyle{plain}
\bibliography{refs}

\end{document}
